\documentclass{article}

% Define the title
\title{DS\_A2\_Q2}
\author{Dongju Ma}
\date{Trimester 1, 2025}

% Load packages
\usepackage{amsmath}
\usepackage{setspace}
\usepackage[a4paper, left=1in, right=1in, top=1in, bottom=1in]{geometry}

% Set global spacing
\setstretch{1.5}
\begin{document}

\maketitle

% Question 2
\section*{2.a}
% Calculate the probability that S_j* contains none of X_1, ..., X_l
Consider the probability that \(S_j^{*}\) contains none of \(X_1, \dots, X_\ell\) is 
\[
    P({X_1, \dots, X_\ell} \cap S_j^{*} = \emptyset)
\] 
and the probability that a single observation in \(S_j^{*}\) is none of \(X_1, \dots, X_\ell\) is
\[
    \frac{n-\ell}{n}
\]
Since we have n observations in \(S_j^{*}\).
Thus, the probability is
\[
    P({X_1, \dots, X_\ell} \cap S_j^{*} = \emptyset) = \left( \frac{n-\ell}{n} \right)^n
\]

\section*{2.b}
% Calculate the expected value of the sample with k observations from X_1, ..., X_l
Since we have \(k\) observations in \(S_j^{*}\) are from \(X_1, \dots, X_\ell\), the sum of these observations is
\[
    \sum_{i=1}^{k} X_i = k E[X_i] = k (\theta - \phi )
\]
And the left \(n-k\) observations are from \(X_{\ell+1}, \dots, X_n\), the sum of these observations is
\[
    \sum_{i=k+1}^{n} X_i = (n-k) E[X_i] = (n-k) \theta
\]
Thus, the sum of all observations is
\[
    \sum_{i=1}^{n} X_i = k (\theta - \phi) + (n-k) \theta = n \theta - k \phi,
\]
which determines that
\[
    \bar{X_j} = \theta - \frac{k \phi} {n}
\]

\section*{2.c}
% Find the value of k makes the estimator \theta unbiased
As \(\bar{X_j}\) is an unbiased estimator of \(\theta\), we have
\[
    E[\bar{X_j}] = \theta
\]
According to 2.b, we have,
\[
    E[\bar{X_j}] = \theta - \frac{k \phi} {n} = \theta \Rightarrow \frac{k \phi}{n} = 0 
\]
which determines that
\[
    k = 0.
\]

\section*{2.d}
% Calculate the expected bias \bar{X_j} as an estimator of \theta
According to the defination of expected bias and 2.b, we have
\begin{align*}
    \text{Bias}(\bar{X_j}) &= E[\bar{X_j}] - \theta \\
    &= \theta - \frac{k \phi} {n} - \theta \\
    &= - \frac{k \phi} {n}
\end{align*}

\end{document}